\section{Identification of the covariance and the Gaussian limit}
\label{sec:gaussian-limit}

Fix a finite observation list and absorb the coefficients of any desired
real linear combination into its tests \(\phi_j\).  We first prove the
Gaussian limit for the single pasted dynamics supplied by
\Cref{thm:marked-expansion,thm:sharp-cumulants}.  Its parameters
\(\Gamma,L\) do not depend on the observation list or a moment order;
only auxiliary complexity and initial-tail truncations in the expansion
may depend on the requested order and accuracy.  The quantitative
comparison with the true flow is proved in \Cref{sec:cutoff}, and
\Cref{sec:transfer} then transfers the exact centering.

\subsection{Records and the limiting one- and two-root measures}

Give every particle a scalar mark, initially \(y=0\), and at time \(t_j\)
perform the update
\begin{equation}\label{eq:limit-mark-update}
 (x,v,y)\longmapsto \Theta_j(x,v,y)
       =(x,v,y+\phi_j(x,v)).
\end{equation}
Collisions preserve the mark of each individual particle.  If several
observations have the same time, their updates may equivalently be
combined.  Write
\[
 M_{\mathrm{obs}}=\sum_j\|\phi_j\|_\infty,
 \qquad \xi=(x,v,y),\qquad |y|\le M_{\mathrm{obs}}.
\]
The construction is only a record of the deterministic trajectories.
It introduces no additional randomness.  On either dynamics,
\begin{equation}\label{eq:limit-record-identity}
        \sum_{i=1}^N y_i(T)
            =\sum_j\sum_{i=1}^N\phi_j(z_i(t_j)).
\end{equation}

Use superscript \(\eps\) for the normalized factorial cumulants of the
pasted marked process.  For every observation-free time, and with the
appropriate one-sided convention at an update,
\Cref{thm:sharp-cumulants} provides limits
\begin{equation}\label{eq:limit-F-J}
               c_1^\eps\longrightarrow F_t,
           \qquad \mu c_2^\eps\longrightarrow J_t.
\end{equation}
Here \(F_t\) is a one-root measure and \(J_t\) is a symmetric signed
two-root measure.  The convergence is against bounded continuous tests
of the records, and the limiting tree measures have the Gaussian
variation bounds in \Cref{thm:sharp-cumulants}.
Before the updates at time zero, \Cref{prop:initial} gives
\begin{equation}\label{eq:limit-initial-F-J}
                 F_0=f_0(x,v)\,\dd x\,\dd v\,\delta_0(\dd y),
                 \qquad J_0=0.
\end{equation}
Indeed \(\mu\|c_{2,0}^{\eps}\|_1=O(\eps)\).

Forgetting records in the finite marked hierarchy commutes with its
one-root regroupings.  By
\Cref{thm:marked-expansion,prop:one-root,lem:sharp-tree-completeness},
the \((x,v)\) marginal of \(F_t\) is \(f_t(x,v)\,\dd x\,\dd v\).
This statement applies also inside a slab: the one-root block map on a
partial slab has the approximation \eqref{eq:dhm-partial-slab}, with
its contact times restricted to that partial slab.  The unmarked dependent
one-root tree terms vanish by \Cref{lem:sharp-tree-completeness}, so
the remaining marginal is precisely that one-root Boltzmann expansion.

Since the first measure in \eqref{eq:limit-F-J} is positive, \(F_t\)
is positive.  Its disintegration supplied by
\Cref{lem:sharp-spatial-traces} is written
\begin{equation}\label{eq:limit-F-disintegration}
                         F_t=\dd x\,F_t(x,\dd v\,\dd y).
\end{equation}
Its unmarked marginal equals \(f_t(x,v)\,\dd v\) at every \(x\), in
the weak continuous version of that disintegration.  In fact equality
holds for almost every \(x\) by the equality of the global marginals,
and both sides are weakly continuous in \(x\).  The same reasoning
allows a positive version of the conditional measure.

We specify the collision products that occur below.  If a factor is
\(F_t\), or \(J_t\) with one root position designated, use its conditional
measure at that position from \Cref{lem:sharp-spatial-traces}.
Multiply two such separate component measures at their common position,
and integrate that position against \(\dd x\).  Include the incoming
flux \(B(v-v_*,\omega)\,\dd\omega\) when the incident variables collide.
These are absolutely defined products.  A uniform conditional Gaussian
bound for one component, the integrated Gaussian bound for the other,
and \(B\le |v|+|v_*|\) give an integrable majorant.  Equivalently, one
can first replace equality of the two positions by a smooth approximate
identity and then take its width to zero.  This characterization and
weak continuity fix the version of the product.  We use no
pointwise product of unspecified almost-everywhere disintegrations.

\subsection{Last-contact evolution}

At a collision write
\[
 \xi=(x,v,y),\quad \xi_*=(x,v_*,y_*),\quad
 \xi'=(x,v',y),\quad \xi_*'=(x,v_*',y_*).
\]
For a marked test \(\Phi\), set
\[
 \Delta\Phi=\Phi(\xi')+\Phi(\xi_*')-\Phi(\xi)-\Phi(\xi_*),
 \qquad D\Phi=v\cdot\nabla_x\Phi,
\]
and define
\begin{equation}\label{eq:limit-marked-H}
 (H_t\Phi)(\xi)
   =\int_{\R^3\times[-M_{\mathrm{obs}},M_{\mathrm{obs}}]\times S^2}
       B(v-v_*,\omega)\Delta\Phi\,
          F_t(x,\dd v_*\,\dd y_*)\,\dd\omega.
\end{equation}
In the next formulas, \(\int BFF_*(\cdots)\) denotes the simultaneous
position integral formed from \eqref{eq:limit-F-disintegration},
including \(\dd x\,\dd\omega\).

\begin{proposition}[Evolution of the marked tree laws]
\label{prop:tree-evolution}
Between observation updates the limiting measures satisfy
\begin{align}
 \partial_tF_t(\Phi)
  &=F_t(D\Phi)+\frac12\int BF_tF_{t,*}\Delta\Phi,
       \label{eq:limit-F-evolution}\\
 \partial_tJ_t(\Phi\otimes\Psi)
  &=J_t\bigl((D+H_t)\Phi\otimes\Psi
                     +\Phi\otimes(D+H_t)\Psi\bigr)\notag\\
  &\quad+\int BF_tF_{t,*}
          \bigl(\Phi(\xi')\Psi(\xi_*')
                         -\Phi(\xi)\Psi(\xi_*)\bigr).
       \label{eq:limit-J-evolution}
\end{align}
The identities hold in transport mild form, and hence in weak
differential form for smooth tests for which the displayed integrals
exist.  At an update \(t_j\),
\begin{equation}\label{eq:limit-pushforwards}
       F_{t_j+}=(\Theta_j)_\#F_{t_j-},
       \qquad
       J_{t_j+}=(\Theta_j\otimes\Theta_j)_\#J_{t_j-}.
\end{equation}
\end{proposition}

\begin{proof}
We derive the identities from the tree integrals, after the limit,
rather than taking a microscopic contact trace.  Start with smooth
bounded marked tests, bounded on the allowed mark interval.  Expand the
independent one-root factors of the highest slab by their definition
for that slab, retaining the canonical lower inputs, and then truncate
the resulting trees to finitely many vertices and bounded total birth energy.
The size restriction thus includes the newly exposed lines, as in
\Cref{lem:sharp-tree-splitting}.  Away from the finitely many
observation times, the test at the upper endpoint depends on current
root phase points and already accumulated marks.

Vary that upper endpoint, or equivalently integrate in the time of the
last contact of the highest slab.  The scattering and non-scattering
choices at a last contact that meets no tested root cancel: no later
contact or update uses either line, so their tests are unchanged, while
the two choices have the same incoming flux and opposite signs.
The change when a tested root is incident is its scattered value minus
its free value.  Deleting that last edge splits a global tree into two
components.  By \Cref{lem:sharp-tree-splitting}, these are exactly the
two component measures before contact, multiplied at the common
position in the sense just specified.  This lemma also supplies the
inverse joining operation and its fresh-label coefficients.
At a finite total-size truncation, the pair of component histories
retains the inherited restriction on their combined size.  The
unrestricted product of the two component measures appears only
after absolute summability removes that auxiliary restriction.

For clarity, the root and label counts can be checked as follows.
\begin{enumerate}[label=\textup{(\alph*)}]
\item With one original root, its partner has been chosen from the
fresh labels; selecting it changes \(1/j!\) into \(1/(j-1)!\).
Dividing the remaining labels between the two component trees uses
\[
 \frac1{(j_1+j_2)!}\binom{j_1+j_2}{j_1}
                         =\frac1{j_1!j_2!}
\]
in each slab.  Thus the contact contribution is
\(\int BFF_*(\Phi(\xi')-\Phi(\xi))\).  Particle interchange,
including \(\omega\mapsto-\omega\), symmetrizes this into the
collision term in \eqref{eq:limit-F-evolution}.
\item With two original roots, a contact between the original pair
is one unordered pair and introduces no fresh-label factor.
Deleting it gives two one-root means.  Its action on the product
test is precisely the last integral in
\eqref{eq:limit-J-evolution}.  There is no factor \(1/2\) in this
unsymmetrized expression for the two designated roots.
\item Suppose the contact is on the first original root and a
new partner.  The second original root lies in one of the two
components obtained by deleting the contact.  If it lies with the
first root, the contribution is \(J\) on that component and \(F\)
on the partner component, with change
\(\Phi(\xi')-\Phi(\xi)\).  If it lies with the new partner, interchange
the two incident variables; the \(J\) component then carries the
change of the test on the other incident line,
\(\Phi(\xi_*')-\Phi(\xi_*)\).  Adding the two placements gives
\(J(H_t\Phi\otimes\Psi)\).  The same argument at the second
original root gives \(J(\Phi\otimes H_t\Psi)\).
\end{enumerate}
If the temporarily designated new partner has no record in one of
these integrands, its component is simply projected on the needed
variables.  Conversely, when its mark occurs after an interchange,
the full recorded component is used.  The projection identity in
\Cref{thm:marked-expansion,lem:sharp-tree-completeness,lem:sharp-tree-splitting}
identifies these as projections of the same \(F,J\); it does not
postulate factorization of the microscopic law at an interface.

Free flight gives the \(D\) terms.  One can formulate the calculation
without differentiating a sum: for an interval \([a,t]\) contained
in one slab and containing no update, freely transport the final
test from \(t\) to the last-contact time \(s\).  The no-contact term
is the measure at \(a\), tested against this free transport, and the
last-contact terms are the integrals of the expressions above over
\(a<s<t\).  For \(J\) use free transport in both test slots.
These are exactly the transport mild versions of
\eqref{eq:limit-F-evolution} and \eqref{eq:limit-J-evolution}.

The finite re-expansion, deletion, and joining were bijections of
labelled trees with their factorial weights.  Passing to the full
sums and the unbounded velocity domain is justified by
\Cref{lem:sharp-summability,lem:sharp-spatial-traces,lem:sharp-tree-splitting}.
In particular the last-contact products are integrable in time and
space with a fixed Gaussian exponent allowance.  The same bounds
justify locally uniform approximation of continuous tests with
smaller Gaussian growth.  They also extend the identities to bounded
tests such as \(y\), by cutting off space and velocity first and using
the integrated variation bounds.

At a slab interface with no contact at its endpoint, the empty upper
slab has the input hierarchy as its no-contact term.  The finite
marked expansion has this identity before regrouping, and the
absolutely convergent tree sums preserve it.  Hence there is no jump
at an artificial interface, and the mild identities concatenate
across all slabs.  Actual observation updates are the deterministic
pushforwards of the finite marked hierarchy; bounded continuous
record convergence gives \eqref{eq:limit-pushforwards}.
\end{proof}

\subsection{Collision normalization and the covariance source}

The pre/postcollision substitution, combined with
\(\omega\mapsto-\omega\), preserves
\(B(v-v_*,\omega)\,\dd v\,\dd v_*\,\dd\omega\).
Particle interchange preserves it with the same change of normal.
Consequently the normalization in \eqref{eq:boltzmann} gives
\begin{equation}\label{eq:limit-weak-normalization}
 \int Q(f,f)\phi=\frac12\int Bff_*\Delta\phi,\qquad
 \int[Q(g,f)+Q(f,g)]\phi=\int g\,K_t\phi
\end{equation}
when the background is \(f=f_t\).  Thus the collision part of the
linearized weak drift has no extra factor \(1/2\).

Restore the same-particle term by defining
\begin{equation}\label{eq:limit-G-definition}
              G_t(\Phi,\Psi)=J_t(\Phi\otimes\Psi)+F_t(\Phi\Psi).
\end{equation}
Then \Cref{prop:tree-evolution} gives
\begin{equation}\label{eq:limit-G-evolution}
 \begin{split}
 \partial_tG_t(\Phi,\Psi)
   ={}&G_t((D+H_t)\Phi,\Psi)+G_t(\Phi,(D+H_t)\Psi)\\
     &+\frac12\int BF_tF_{t,*}\Delta\Phi\,\Delta\Psi.
 \end{split}
\end{equation}
Here is a direct verification of the coefficient and sign of its
source.  Write \(P=\Phi+\Phi_*\) and \(R=\Psi+\Psi_*\).
The collision change of the diagonal products from the \(F\) equation,
plus the source in the \(J\) equation, symmetrized by particle
interchange, is
\[
                 \frac12\int BFF_*\,\Delta(PR).
\]
The diagonal contributions in the two \(H_t\) drift terms are
\[
                 \frac12\int BFF_*
                       \bigl(R\Delta P+P\Delta R\bigr).
\]
Their difference is the last term of
\eqref{eq:limit-G-evolution}, because
\(\Delta(PR)=R\Delta P+P\Delta R+\Delta P\,\Delta R\).
This uses neither detailed balance nor an equilibrium identity.
The source is positive semidefinite, since \(F_t\) is positive.
At time zero, before updates,
\begin{equation}\label{eq:limit-G-initial}
                   G_0(\Phi,\Psi)=\int f_0(z)\Phi(z,0)\Psi(z,0)\,\dd z,
\end{equation}
and at an update \(G\) pushes forward in its two test slots.

\subsection{Backward pulses and all observation times}

We only need \eqref{eq:limit-G-evolution} on the spaces
\[
 \mathcal V_a=\{\,\Phi(x,v,y)=c\,y+h(x,v):
                   c\in\R,\ h\in E_a^0\,\},
 \qquad \|\Phi\|_{\mathcal V_a}=|c|+\|h\|_a.
\]
Let \(\alpha>0\) be a common exponent in the integrated and conditional
Gaussian variation bounds for \(F,J\).  Choose \(0<a<b<\beta\)
with \(2b<\alpha\), leaving a positive exponent margin for the
collision flux.  The factor two also controls the diagonal
term \(F_t(hk)\), where both tests use the same velocity.
The mark is constant on each particle between updates, and its pair
sum is conserved at a collision.  The pointwise marginal statement
following \eqref{eq:limit-F-disintegration} therefore implies
\begin{equation}\label{eq:limit-restricted-drift}
 (D+H_t)(c\,y+h)=A_th,\qquad
 \frac12\int BF_tF_{t,*}\Delta(c\,y+h)\Delta(d\,y+k)=C_t(h,k).
\end{equation}
At the pulse \(t_j\), backward substitution changes \(c\,y+h\)
to \(c\,y+h+c\phi_j\).

We justify using the backward propagator in this possibly nonsmooth
test class.  The Gaussian variation bounds for \(F,J\) make \(G_t\)
a uniformly bounded bilinear form on \(\mathcal V_b\).
Its transport mild identity extends from the smooth
tests in \Cref{prop:tree-evolution} by local uniform approximation
with a small Gaussian exponent allowance.  On the unmarked parts its
collision operator is \(K_t\), so the scale estimate
\eqref{eq:target-K-estimate} applies.
Iterating the bilinear mild identity amounts to choosing one of the
two slots for each collision factor.  An \(n\)-fold remainder is
bounded by a constant times
\[
             \frac{(CT)^n}{n!}
                  \bigl(1+\sqrt{n/\delta}\bigr)^n
\]
for the fixed total exponent allowance \(\delta=b-a>0\); the \(2^n\)
slot choices are included in \(C^n\).  It tends to zero.
The source iterations have the same summable bound.
Thus dual propagation of \eqref{eq:limit-G-evolution} is justified
by its convergent mild series.  In particular it needs no spatial
or velocity differentiability of \(U(r,t)\phi\).

Starting with the test \(y\) after the final observation, the backward
test at a time \(r\) away from the pulses is \(y+h_r\), where
\begin{equation}\label{eq:limit-backward-pulse-test}
                   h_r=\sum_{j:t_j\ge r}U(r,t_j)\phi_j.
\end{equation}
The same convention at \(r=0\) includes observations made at time zero.
The exact values assigned at positive pulse times do not change a
time integral.  Combining the preceding dual propagation, the
initial value \eqref{eq:limit-G-initial}, and the pulse substitutions
gives
\begin{equation}\label{eq:limit-variance-formula}
                  G_T(y,y)
                  =\int f_0h_0^2
                         +\int_0^T C_r(h_r,h_r)\,\dd r.
\end{equation}
After the last pulse the unmarked backward part is zero, so one may
equally use the last observation time in place of \(T\).

By the explicit Gaussian construction
\eqref{eq:target-isonormal-formula}, the random variable
\(\sum_j\zeta_{t_j}(\phi_j)\) is
\[
                    I_0(h_0)+I(\Delta h_r).
\]
Its variance is precisely \eqref{eq:limit-variance-formula}.
This identifies correlations between different observation times,
as well as observations made repeatedly on the same particle.

\subsection{Ordinary cumulants and Gaussian convergence}

\begin{proposition}[Gaussian limit on the pasted dynamics]
\label{prop:cumulant-clt}
For the single pasted dynamics fixed above and every fixed observation
list,
\[
 \frac{\sum_i y_i(T)-\E\sum_i y_i(T)}{\sqrt\mu}
       \ \Longrightarrow\
           \sum_j\zeta_{t_j}(\phi_j).
\]
All centered moments of the scalar random variable on the left
converge to the corresponding Gaussian moments.
\end{proposition}

\begin{proof}
Put \(S^\eps=\sum_i y_i(T)\).  The bound
\(|S^\eps|\le M_{\mathrm{obs}}N\) and the partition function imply
that \(S^\eps\) has exponential moments for each fixed \(\eps\).
For the normalized factorial cumulants at the final marked time,
the exact factorial-to-ordinary identity is
\begin{equation}\label{eq:limit-factorial-ordinary}
 \kappa_k(S^\eps)=
   \sum_{\pi\in\mathcal P([k])}\mu^{|\pi|}
       \int c_{|\pi|}^{\eps}(\dd\xi_1\cdots\dd\xi_{|\pi|})
                     \prod_{B\in\pi} y_B^{|B|}.
\end{equation}
Here one associates the blocks of \(\pi\), in any fixed canonical
order, with the distinct particle slots in the integral.
For a direct derivation, write
\(\exp(\lambda S^\eps)=\prod_i[1+(\exp(\lambda y_i)-1)]\),
insert the factorial generating series, and take its logarithm.
Each factor \(\exp(\lambda y)-1\) must receive a nonempty set of the
\(k\) differentiations at zero.  Distributing those sets gives
exactly the partitions in \eqref{eq:limit-factorial-ordinary};
their ordering cancels the factorial denominator of the factorial
series.  This is also a finite formal identity through degree \(k\),
so no uniform radius of convergence in \(\eps\) is being assumed.

All the powers of \(y\) in \eqref{eq:limit-factorial-ordinary} are
bounded on the common mark interval; they can be extended to bounded
continuous functions on the full mark line.  By
\Cref{thm:sharp-cumulants}, every term on its right is \(O_k(\mu)\).
For the centered normalized variable
\[
                     X^\eps=(S^\eps-\E S^\eps)/\sqrt\mu
\]
the first cumulant is zero.  The second is
\[
 \kappa_2(X^\eps)
       =c_1^\eps(y^2)+\mu c_2^\eps(y\otimes y)
       \longrightarrow F_T(y^2)+J_T(y\otimes y)=G_T(y,y).
\]
For \(k>2\),
\begin{equation}\label{eq:limit-higher-cumulants}
                  \kappa_k(X^\eps)=O_k(\mu^{1-k/2})\longrightarrow0.
\end{equation}
The statement that these estimates hold on one dynamics is material:
the bases determining \(\Gamma,L\) in
\Cref{thm:marked-expansion,thm:sharp-cumulants} are independent of
the fixed factorial order.  Choosing a larger auxiliary complexity
threshold or a larger \(D\) only changes the error representation
used to prove \eqref{eq:limit-higher-cumulants}; it does not change
the random variable \(X^\eps\).

The finite moment--cumulant identities now give every centered
Gaussian moment with variance \(\sigma^2=G_T(y,y)\).
Here the usual Gaussian moment-determinacy argument can be made
explicit.  Truncate the characteristic functions of \(X^\eps\)
and of a centered Gaussian \(Z\) with variance \(\sigma^2\) at
degree \(2n-1\).  Taylor's integral remainder on the real axis is
bounded by \(|u x|^{2n}/(2n)!\), so moment convergence gives
\[
 \limsup_{\eps\to0}
  \left|\E e^{iuX^\eps}-\E e^{iuZ}\right|
 \le \frac{2|u|^{2n}\E Z^{2n}}{(2n)!}
 =\frac{2|u|^{2n}\sigma^{2n}}{2^n n!}.
\]
Letting \(n\to\infty\) proves characteristic-function convergence
for every real \(u\), including when \(\sigma=0\).
It implies the asserted weak convergence, with variance identified
by \eqref{eq:limit-variance-formula}.
\end{proof}

The Gaussian limit is now established for the pasted dynamics.  To
obtain the same law for the true hard-sphere flow, it remains to compare
not only their observations but also their exact expectations.  The
next section proves a discrepancy probability smaller than the
\(\eps\) scale; \Cref{sec:transfer} combines that estimate with the
initial particle-number moments.
